% =========================================================
% Proof file: prf-scalar-times-zero-vector.tex
% =========================================================

\newpage
\phantomsection
\label{prf:scalar-times-zero-vector}

\begin{remark*}[Return]
\hyperref[thm:scalar-times-zero-vector]{Return to Theorem}
\end{remark*}

\begin{theorem*}[$a0=0$]
For every $a \in \mathbf{F}$,
\[
a0=0.
\]
\end{theorem*}

\begin{proof}
\textbf{Professional Standard Proof.}~\\
Let $a \in \mathbf{F}$. By distributivity,
\[
a(0+0)=a0+a0.
\]
Because $0+0=0$ in $V$, this becomes
\[
a0=a0+a0.
\]
Add the additive inverse of $a0$ to both sides:
\[
-(a0)+a0=-(a0)+(a0+a0).
\]
By associativity,
\[
0=\bigl(-(a0)+a0\bigr)+a0=0+a0=a0.
\]
Therefore
\[
a0=0.
\]
\end{proof}

\begin{proof}
\textbf{Detailed Learning Proof.}~\\

\textbf{Step 1.} Fix an arbitrary scalar.\\
Let $a \in \mathbf{F}$. The goal is to determine the vector $a0$.

\textbf{Step 2.} Use distributivity over vector addition.\\
The vector space axiom gives
\[
a(u+v)=au+av
\]
for all $u,v \in V$. Apply this with $u=0$ and $v=0$:
\[
a(0+0)=a0+a0.
\]

\textbf{Step 3.} Simplify the vector expression inside the parentheses.\\
Because $0$ is the additive identity in $V$,
\[
0+0=0.
\]
Hence
\[
a0=a0+a0.
\]

\textbf{Step 4.} Add the additive inverse of $a0$ to both sides.\\
Because $a0 \in V$, it has an additive inverse $-(a0)$. Adding this vector to both sides gives
\[
-(a0)+a0=-(a0)+(a0+a0).
\]

\textbf{Step 5.} Simplify both sides.\\
On the left,
\[
-(a0)+a0=0.
\]
On the right, associativity gives
\[
-(a0)+(a0+a0)=\bigl(-(a0)+a0\bigr)+a0=0+a0=a0.
\]
Therefore
\[
0=a0.
\]

\textbf{Step 6.} State the conclusion.\\
Thus every scalar sends the zero vector to the zero vector:
\[
a0=0.
\]
\end{proof}

\begin{remark*}[Proof structure]
The proof is direct. Distributivity yields an equation of the form $w=w+w$ with $w=a0$, and additive-inverse cancellation then forces $w=0$.
\end{remark*}

\begin{remark*}[Dependencies]
\hfill
\begin{itemize}
  \item \hyperref[def:vector-space]{Definition of Vector Space}
  \item \hyperref[thm:unique-additive-inverse]{Unique Additive Inverse}
\end{itemize}
\end{remark*}